\section{四枚被覆の全形式標準形と $S_4$}\label{sec:four-versal-v23}

この節では、選んだ十三変数の基底 $A_0$ 上の有限代数 $\Lambda$ を調べる。
四枚被覆としての構造は、四変数の最小versal変形と九変数の自明な方向に分離できる。
ここで変形とは、特殊ファイバーの同一視を持つ有限平坦可換代数の変形を指す。
各変形に原点の切断を要求しない。また、以下の標準形を与える座標と分類写像は一意ではない。

\subsection{長さ四の完全交叉の最小versal変形}

\begin{lemma}\label{lem:four-miniversal-v23}
標数零の体 $k$ 上で
\[
 D_0=k[[X,Y]]/(XY,Y^2-2X^2)
\]
と置く。$D_0$ の可換代数変形の接空間は四次元である。
次の族は最小versal変形であり、その全空間は四次元正則局所環である。
\begin{align}
 R&=k[[u,v,w,z]],\\
 \mathscr D&=R[[X,Y]]/(XY-u,\ Y^2-2X^2-v-wX-zY).
 \label{eq:four-miniversal-v23}
\end{align}
$\mathscr D$ は $R$ 上有限自由階数四である。
\end{lemma}

\begin{proof}
完全交叉の余接複体を使うと、接空間は
\[
 T^1(D_0)=\coker\left[
 D_0^2\xrightarrow{
 \left(\begin{smallmatrix}Y&X\\-4X&2Y\end{smallmatrix}\right)}D_0^2
 \right].
\]
この余核では $(X,0)=(0,-2Y)$、$(Y,0)=(0,4X)$ であり、
二つの成分の二次項は全て消える。従って
\begin{equation}
 [(1,0)],\quad[(0,1)],\quad[(0,X)],\quad[(0,Y)]
 \label{eq:four-T1-basis-v23}
\end{equation}
が基底となる。この記述は完全交叉の余接複体の標準的な二項表示
\cite[Tag 08SH]{Stacks}による。

最小versalityを、この場合には直接確認できる。
完備Noether局所 $k$-代数 $C$ 上の有限平坦変形を考える。
$X,Y$ を持ち上げると $C[[X,Y]]$ から全射があり、
平坦性とNakayamaの補題により、その核は中心の正則列の持上げ二本で生成される。
従って変形を、持ち上げた $X,Y$ と二つの関係で表示できる。
基底の極大イデアルを $\mathfrak m_C$ とし、標準形を
$\mathfrak m_C^n$ を法として構成したとする。
次の次数の二つの誤差を $D_0^2\otimes_k
\mathfrak m_C^n/\mathfrak m_C^{n+1}$ で考え、
その $T^1$ 成分を四つの係数 $u,v,w,z$ の変更で吸収する。
残りのJacobian像は $X,Y$ の変更で消え、
二つの中心関係の倍は可逆な関係基底の変更で消える。
この操作を帰納的に続けると完備性により全形式の標準形を得る。
接写像が基底\eqref{eq:four-T1-basis-v23}を与えるので最小である。
これは一意な分類写像の主張ではない。

最後に $u,v$ を消去すると
$\mathscr D\simeq k[[X,Y,w,z]]$ である。
閉ファイバーの長さが四なので、完備Nakayamaにより有限であり、
両環の正則性と次元一致から平坦となる。
局所環上の有限平坦加群は自由で、その階数は閉ファイバーの長さ四である。
\end{proof}

\subsection{実際の族のKodaira--Spencer写像}

\begin{proposition}\label{prop:actual-four-KS-v23}
実際の $\Lambda/A_0$ のKodaira--Spencer写像は全射である。
特に $t_8,t_{11},t_{13},t_{14}$ の四方向だけで $T^1(D_0)$ を生成する。
\end{proposition}

\begin{proof}
この有限線形計算は、元の十六変数族の関係式から行う。
$a=t_5,b=t_{15},c=t_{16}$ と書き、閉ファイバーでの五式を
\[
 G=(F_1,F_2,F_{1x},F_{1y},F_{2x}),\qquad
 \mathbf q=(a,b,c,x,y)
\]
とする。$D_0$ 上では
$a=11x^2,b=x^2,c=0$ である。
五変数五式の余接表示から
\[
 T^1(D_0)=\coker\bigl(D_0^5\xrightarrow{J_{\mathbf q}G}D_0^5\bigr).
\]
基底 $(1,x,y,x^2)$ を使うと $J_{\mathbf q}G$ は二十次正方行列で、
有理数上の階数は十六である。各基底方向 $t_i$ の像は
$[\partial G/\partial t_i]$ で与えられる。
左零化空間のある基底で、上記四方向の列は
\[
 \begin{pmatrix}
 0&1&0&0\\
 1/5&1&-1/5&1/5\\
 -44/25&-51/5&-6/25&-44/25\\
 0&26/5&-4&-4
 \end{pmatrix},
 \qquad \det=-\frac85.
\]
従って接写像は全射である。

検証器\nolinkurl{four_cover_versal_v23.py}は、
実際の混合関係係数から二十次Jacobian、左零化空間、十三本の接方向を再構成し、
Fraction演算と独立な有理行列演算で階数と小行列式を検査する。
入力の高次打切りはこの計算に影響しない。
閉ファイバーで $(a,b,c)^2=0$、$(x,y)^3=0$ なので、
一回のパラメータ微分と高々二回の $x,y$ 微分の後に残り得る項は、
保存された実際の混合jetに全て含まれる。
\end{proof}

\begin{theorem}[四枚被覆の標準形]\label{thm:four-normal-form-v23}
ある形式座標と $A_0$-代数同型により
\begin{align}
 A_0&\simeq\mathbb C[[u,v,w,z,s_1,\ldots,s_9]],\\
 \Lambda&\simeq
 \frac{A_0[[X,Y]]}{(XY-u,\ Y^2-2X^2-v-wX-zY)}.
 \label{eq:actual-four-normal-form-v23}
\end{align}
従って $\Lambda/A_0$ は、式\eqref{eq:four-miniversal-v23}の族に
九変数の形式円板を加えたものである。
\end{theorem}

\begin{proof}
補題\ref{lem:four-miniversal-v23}により、分類写像 $R\to A_0$ と
$\Lambda\simeq\mathscr D\widehat\otimes_R A_0$ を取れる。
命題\ref{prop:actual-four-KS-v23}と $R$ の最小性から、
$u,v,w,z$ の像の一次項は $\mathfrak m_{A_0}/\mathfrak m_{A_0}^2$ で独立である。
これらを十三個の正則パラメータへ延ばし、形式的逆関数定理を適用すると
$A_0\simeq R[[s_1,\ldots,s_9]]$ を得る。
これは $\Lambda$ のみに関する標準形であり、
元の十一点族 $T/B$ が相対potentialを持つことは意味しない。
\end{proof}

\subsection{ramification divisorと一般Galois群}

\begin{corollary}\label{cor:four-ramification-v23}
$\Spec\Lambda\to\Spec A_0$ のramification divisorは、形式的に
三次元の普通二重点と九次元形式円板の積である。
具体的には $\Lambda\simeq\mathbb C[[X,Y,w,z,s_1,\ldots,s_9]]$ の中で
\[
 4X^2+2Y^2+wX-zY=0
 \label{eq:four-ramification-v23}
\]
と書ける。その特異部分は $X=Y=w=z=0$ で九次元である。
\end{corollary}

\begin{proof}
有限写像は $u=XY$、$v=Y^2-2X^2-wX-zY$ である。
その $X,Y$ に関するJacobian行列式は
$2Y^2-zY+4X^2+wX$ である。
ここで線形座標
\[
 \alpha=X,\qquad\beta=4X+w,\qquad
 \gamma=Y,\qquad\delta=2Y-z
\]
を取ると、式は $\alpha\beta+\gamma\delta=0$ となる。
従って三次元の普通二重点と九つの自明な方向の積である。
ここで記述しているのは被覆の全空間内のramification divisorであり、
基底内のbranch divisorそのものの標準形ではない。
\end{proof}

\begin{theorem}\label{thm:four-S4-v23}
四枚被覆 $\Lambda/A_0$ の一般ファイバーのGalois閉包の群は $S_4$ である。
\end{theorem}

\begin{proof}
まず九つの自明な変数を除いた族を考える。
$u\ne0$ では $Y=u/X$ と消去でき、$X$ の方程式は
\[
 2X^4+wX^3+vX^2+zuX-u^2=0.
 \label{eq:four-quartic-v23}
\]
完備化によってGalois群が小さくなる可能性も除くため、
順序付きの四根 $r_1,\ldots,r_4$ の代数を直接調べる。
$e_i$ を四根の基本対称式とすると、Vietaの式は
\[
 e_1=-w/2,\qquad e_2=v/2,\qquad
 2e_3+uz=0,\qquad 2e_4+u^2=0.
\]
従って $w,v$ を消去した多項式環での分解代数は
\[
 C=\frac{\mathbb C[r_1,r_2,r_3,r_4,u,z]}
 {(u^2+2e_4,\ uz+2e_3)}.
\]
最初の式は既約であり、二つの式は正則列をなす。
一方 $u=0$ のファイバーでは $e_3=e_4=0$ となる。
これは少なくとも二つの根が零という条件であり、$z$ を加えて次元三である。
$C$ は次元四の完全交叉なので、$u$ は非零因子である。
また
\[
 C[u^{-1}]
 \simeq
 \frac{\mathbb C[r_1,r_2,r_3,r_4,u,u^{-1}]}
 {(u^2+2r_1r_2r_3r_4)}
\]
は整域である。実際、四根の積は分数体内の平方ではない。
従って $C$ 自身が整域である。

$r_i,z$ の重みを一、$u$ の重みを二とすると $C$ は正次数付き整域である。
重みの位相は原点の極大イデアル進位相と同値であり、
その完備化 $\widehat C$ の随伴次数環は $C$ である。
従って $\widehat C$ も整域となる。
四次のモニック多項式の順序付き分解代数は基底上自由階数 $4!=24$ である
（根を一つずつ加え、残りのモニック因子で割る構成による）。
特殊ファイバーの全ての根は零なので、$\widehat C$ は
$\mathbb C[[u,v,w,z]]$ 上のその分解代数であり、階数二十四を持つ。
判別式が零多項式でないことは、$u=1,w=v=z=0$ で $2X^4-1$ となることから分かる。
この非零性は整域の完備化でも保たれるので、一般の四根は相異なり、
$S_4$ はそれらを忠実に置換する。
分数体の次数も二十四なのでGalois群は正確に $S_4$ である。

九つの自明な形式変数を加えた順序付き分解代数は
$\widehat C[[s_1,\ldots,s_9]]$ で、引き続き整域である。
したがって $\Lambda/A_0$ に対しても同じ結論となる。
\end{proof}

この $S_4$ は、十三変数への投影で現れる四つのcorank二点のmonodromyである。
元の残留十一点の $S_{11}$、および全二十一点の $S_{21}$ とは対象が異なる。
